\section{Conclusion}

The code conductor turns the transversal-group problem into a
one-dimensional linear test.  For the present half-rate MDS code,
\(\Cond(C,C^\perp)=(C^{\star2})^\perp\) is either zero or a full-support line.
Equivalently, the coordinatewise endomorphism algebra of the CSS label space
is either \(\F_q\times\F_q\) or \(M_2(\F_q)\).  Their determinant-one units
give the split torus and \(\mathrm{SL}_2(q)\), respectively.  Stabilizer-AME
rigidity then proves that these Clifford constructions exhaust every
product-unitary logical implementation.

For six coordinates, the same dichotomy becomes the conic criterion in the
six-arc model.  Further geometric and symmetry refinements show how the
all-length criterion organizes finer finite geometry without depending on
those calculations.  Fixed-copy scalar
contractions are generically constant, whereas the operator-valued reduced
states retain the local Weyl directions.  This explains why the rigidity
argument succeeds where bounded-degree scalar invariants do not.

Over prime-power fields the same conductor test sees every additive label
map.  Additive conductors decompose into Galois sectors, each block is
semilinear over the cross-ratio field \(K_C\), and the fixed-coordinate
label group is \(\mathrm{GL}_d(K_C)\), \(\mathrm{Sp}_{2d}(K_C)\), or
\(\mathrm{GU}_{2d}(K_0)\).  New symmetries beyond the \(\F_q\)-linear ones
come only from subfield structure or from the Hermitian relation
\(C^\perp\sim\tau C\).  When \(K_C=\F_q\), all but \(2(q-1)\) elements of
the Hermitian group \(\mathrm{GU}_2(\sqrt q)\) mix Galois sectors.

Several questions remain.  For coordinate-permuting symmetries, the nonabelian factor set gives an
exact splitting criterion, and the integral \(H_3\) family now has a uniform
complement; a conceptual classification beyond that family is open.  Beyond
MDS codes, the same coordinatewise algebra may
acquire a square-zero off-diagonal radical, suggesting Borel subgroups as the
first intermediate field-linear CSS symmetry type.  A general theorem must
also account for decomposable codes, positive logical dimension, and
prime-power trace-symplectic effects without the MDS property, and is left
for future work.  For the pencil over extension fields, it remains open
whether every Galois match of \(z\) is realized by a local Clifford.
